% Supplementary material of main.tex — generic template (plain article class). To move to a venue class
% (IEEEtran, acmart, llncs, ...), keep the body and replace the preamble and the front matter.
% Build:  latexmk -pdf main.tex
% Figures: figures/paper_figures.ipynb (style: figures/paper_style.py; dataset_construction.py runs on the server)
% Notes for the authors would be in \todo{...} (red). The previous draft is in _da_cancellare/.
\documentclass[11pt,a4paper]{article}
\usepackage[margin=2.5cm]{geometry}
\usepackage[T1]{fontenc}
\usepackage[utf8]{inputenc}
\usepackage{lmodern}
\usepackage{amsmath,amssymb}
\usepackage{graphicx}
\usepackage{booktabs}
\usepackage{xcolor}
\usepackage{tikz}
\usetikzlibrary{arrows.meta,positioning,shapes.misc}
\usepackage[numbers,sort&compress]{natbib}
\usepackage{xr-hyper}
\usepackage[colorlinks=true,linkcolor=blue!50!black,citecolor=blue!50!black,urlcolor=blue!50!black]{hyperref}
\def\UrlBreaks{\do\/\do-\do\_\do.}
\usepackage{authblk}
\usepackage{orcidlink}
\renewcommand\Affilfont{\small}

\externaldocument[M-]{main}
\newcommand{\todo}[1]{\textcolor{red}{[\textbf{TODO:} #1]}}
\newcommand{\anon}{\operatorname{anon}}

\renewcommand{\thesection}{S\arabic{section}}
\renewcommand{\thetable}{S\arabic{table}}
\renewcommand{\thefigure}{S\arabic{figure}}

\title{Supplementary material for ``Traces of Reliability: Predicting the Rating of Unseen News Publishers from the
Anonymized Text of Their Articles''}
\author{}
\date{}

\begin{document}
\maketitle



\section{More metrics}\label{app:measures}
Table~\ref{tab:more} adds three metrics to Table~\ref{M-tab:main}, on the same outlets, predictions and resamples: the
coefficient of determination $R^2$, the root mean squared error (RMSE), which weighs large errors more than the MAE, and
the area under the ROC curve (AUC) for separating the reliable outlets (score $\geq 60$, the NewsGuard threshold, 65\% of
the outlets) from the others. The AUC uses only the order of the predictions, so it is not affected by their
compression. With $R^2$ and the RMSE the models are in the same order as with the MAE. On the binary distinction the
large frozen models are only slightly ahead of TF-IDF (0.876 and 0.880 against 0.860): their advantage lies more in how
close the predicted score is to the rating than in the side of 60 on which an outlet falls.

\begin{table}[h]
\centering
\caption{More metrics per outlet. Each row is a model, as in Table~\ref{M-tab:main}; brackets give 95\% bootstrap
intervals. Higher $R^2$ and AUC and lower RMSE are better.}\label{tab:more}
\resizebox{\linewidth}{!}{%
\begin{tabular}{lccc}
\toprule
model & $R^2$ & RMSE (points) & AUC, score $\geq 60$ \\
\midrule
no text (training mean) & 0.000 [$-$0.023, 0.000] & 27.0 [24.8, 29.3] & 0.501 [0.422, 0.577] \\
no text (training median) & $-$0.058 [$-$0.130, $-$0.013] & 27.8 [25.1, 30.5] & 0.500 [0.500, 0.500] \\
orientation only & 0.223 [0.101, 0.327] & 23.8 [21.3, 26.5] & 0.688 [0.611, 0.762] \\
Mistral-Small-24B-Instruct, zero-shot & $-$0.085 [$-$0.261, 0.031] & 28.2 [26.3, 30.0] & 0.665 [0.589, 0.738] \\
\midrule
TF-IDF + ridge & 0.339 [0.273, 0.392] & 22.0 [20.1, 24.0] & 0.860 [0.807, 0.907] \\
ModernBERT (fine-tuned) & 0.341 [0.253, 0.415] & 22.0 [19.9, 24.0] & 0.835 [0.781, 0.884] \\
ModernBERT (frozen) + ridge & 0.308 [0.236, 0.365] & 22.5 [20.6, 24.5] & 0.824 [0.769, 0.875] \\
Qwen3-Embedding-0.6B (frozen) + ridge & 0.268 [0.209, 0.316] & 23.1 [21.1, 25.1] & 0.838 [0.782, 0.891] \\
Qwen3-Embedding-8B (frozen) + ridge & 0.437 [0.368, 0.496] & 20.3 [18.4, 22.2] & 0.880 [0.831, 0.923] \\
Mistral-Small-24B (frozen) + ridge & 0.475 [0.406, 0.533] & 19.6 [17.7, 21.5] & 0.876 [0.827, 0.920] \\
\bottomrule
\end{tabular}}
\end{table}

\section{Statistical tests}\label{app:stats}
Intervals are 95\% percentile bootstrap intervals over outlets, with 10,000 resamples for the results with all the
articles and 2,000 for the curves by number of articles. Two models are compared on the same resampled outlets (paired
bootstrap), with a two-sided bootstrap $p$-value. The $p$-values are corrected with the Holm method within three
families: the comparisons of every model with TF-IDF on the four metrics of Table~\ref{M-tab:main} (24 tests), the
curves of each model against TF-IDF (16 tests per model), and the comparisons behind specific statements of the text
(24 tests). A comparison that belongs to two families, such as a model against TF-IDF with all the articles, is
reported from the first. Our primary comparison is the MAE of the large frozen models against TF-IDF. The uncorrected
intervals of Table~\ref{M-tab:ladder} and Section~\ref{app:anon} are stated as such.

The comparisons of Section~\ref{M-sec:stage2} and Section~\ref{M-sec:k} in detail:
\begin{itemize}
  \item \textbf{Mistral-Small-24B against TF-IDF, all articles:} MAE $-2.0$ [$-2.9$, $-1.2$]; $\rho$ $+0.060$ [0.020,
        0.104] and share within 20 points $+6.6$, not significant after correction (Holm $p = 0.07$ and $0.18$).
  \item \textbf{Qwen3-Embedding-8B against TF-IDF, all articles:} MAE $-1.5$ [$-2.3$, $-0.7$] and share within 20 points
        $+8.4$, both significant; $\rho$ $+0.037$, not significant.
  \item \textbf{Qwen3-Embedding-8B against Qwen3-Embedding-0.6B:} MAE $-2.5$ [$-3.2$, $-1.7$] and $\rho$ $+0.085$ [0.039,
        0.136], both significant. Mistral-Small-24B is not significantly better than Qwen3-Embedding-8B, and the frozen
        ModernBERT is not worse than Qwen3-Embedding-0.6B.
  \item \textbf{Frozen Mistral-Small-24B against its zero-shot version:} $+0.39$ in $\rho$ and $-8.6$ points of error.
  \item \textbf{By number of articles:} with five articles per outlet, Mistral-Small-24B has a lower error than TF-IDF
        with all the articles ($-1.7$ [$-2.6$, $-0.9$]) and a higher $\rho$ (0.71 against 0.67, not significant).
        Qwen3-Embedding-8B is significantly better than TF-IDF in ranking up to ten articles, and Qwen3-Embedding-0.6B has
        a significantly larger error than TF-IDF from two articles on.
\end{itemize}

\section{The anonymization in detail}\label{app:anonfig}
The three passes of Section~\ref{M-sec:anon} work as follows:
\begin{enumerate}
  \item links and e-mail addresses are found with regular expressions;
  \item the outlet's own name is derived from its domain and matched written joined or spaced, with or without ``the''; a
        name that is also a common English word (e.g.\ \emph{The Week}, \emph{Time}) is replaced only when capitalized,
        unless the outlet itself writes it in lowercase;
  \item GLiNER (\texttt{gliner\_large-v2.1}, threshold 0.5) is given the labels \emph{news outlet} and \emph{journalist}
        and applied to windows of 200 words overlapping by 30, since it reads at most 384 tokens; of two overlapping
        spans, the one with the higher score is kept. A span is replaced only if it contains a capitalized word that is
        not a pronoun and does not touch a placeholder of passes 1--2, because the recognizer also tags pronouns such as
        \emph{I}, common nouns such as \emph{reporters} and the placeholders themselves.
\end{enumerate}
Passes 1--2 replace 17,780 mentions of the outlet's own name, 8,559 links and 398 e-mail addresses. GLiNER replaces
41,278 mentions of news outlets and 33,816 of journalists, in 71.8\% of the articles; of the spans it proposes, 73\% of
the outlets and 92\% of the journalists pass the filter. The recognizer touches most articles (other outlets in 54\% and
journalists in 45\%), and the own-name rule 24\% (Figure~\ref{fig:anonymization}). Links are more frequent in the least
reliable outlets (1.4 per 1,000 words in the 0--39 band, 0.3 in the 75--100 bands), while the outlets with a perfect
score cite other outlets and name journalists most often (6.8 placeholders per 1,000 words for 0--39 and 8.8 for 100).

The name of the publishing outlet appears in 37.8\% of the original articles, in 1.7\% after passes 1--2 and in 1.6\%
after the full anonymization, mostly where the name is also an ordinary word. This check uses the same name as the
masking, so it cannot see what the masking misses. An independent check with the brand name recorded by NewsGuard finds
it in 2.1\% of the articles, concentrated in 17 outlets whose brand differs from their domain, such as think tanks and
advocacy organizations named in most of their own articles. Social-media handles (in 9.5\% of the articles) and some
aliases are not masked. The models do not read the titles, where signatures are most frequent.

\begin{figure}[t]
\centering
\includegraphics[width=\linewidth]{figures/anonymization.pdf}
\caption{What the anonymization replaces. (a) Share of articles with at least one replacement of each kind (passes 1--2:
links, e-mail addresses, the outlet's own name; pass 3: other outlets and journalists). (b) Share of articles whose
title or text still contains the name of their own outlet, after each pass. (c) Replacements per 1,000 words, by
NewsGuard score band of the outlet.}\label{fig:anonymization}
\end{figure}

\section{Who is left out by the cleaning}\label{app:selection}
Table~\ref{tab:selection} describes the 1,293 political outlets sent to the scraper by the step at which they left the
dataset. The 226 outlets are not a random sample of them (Mann--Whitney $p < 10^{-25}$ for the score and $\chi^2$
$p < 10^{-14}$ for the orientation). The outlets that could not be collected and those removed with the duplicates are
much less reliable than those in the dataset, and the second group is almost entirely the network of 164 copies of
Section~\ref{M-sec:groups}.

\begin{table}[h]
\centering
\caption{Who is left out by the cleaning. Each row is the step at which outlets left the dataset (the last row is the
dataset); the columns give their number, their mean NewsGuard score and their shares (\%) by level, by orientation and
based in the United States.}\label{tab:selection}
\resizebox{\linewidth}{!}{%
\begin{tabular}{lrrrrrrrrrrr}
\toprule
removed at & outlets & score & $l_1$ & $l_2$ & $l_3$ & $l_4$ & $l_5$ & left & right & none & US \\
\midrule
no article collected & 619 & 41.4 & 43 & 35 & 5 & 13 & 3 & 7 & 50 & 43 & 68 \\
text extraction & 6 & 78.2 & 0 & 33 & 0 & 50 & 17 & 0 & 67 & 33 & 50 \\
English & 91 & 73.6 & 11 & 14 & 9 & 45 & 21 & 20 & 22 & 58 & 15 \\
100--1,500 words & 12 & 55.0 & 25 & 17 & 25 & 25 & 8 & 8 & 67 & 25 & 67 \\
published from 2025 & 34 & 56.4 & 21 & 41 & 6 & 24 & 9 & 29 & 26 & 44 & 76 \\
exact duplicates & 170 & 37.0 & 89 & 2 & 2 & 6 & 1 & 7 & 5 & 88 & 98 \\
near duplicates & 14 & 27.9 & 93 & 0 & 7 & 0 & 0 & 0 & 29 & 71 & 100 \\
at least 10 articles & 60 & 57.5 & 22 & 35 & 10 & 25 & 8 & 27 & 35 & 38 & 77 \\
one outlet per group & 61 & 47.8 & 13 & 69 & 7 & 10 & 2 & 36 & 59 & 5 & 97 \\
\textbf{in the dataset} & 226 & 68.5 & 15 & 19 & 13 & 35 & 17 & 30 & 44 & 26 & 79 \\
\bottomrule
\end{tabular}}
\end{table}

\section{Implementation details}\label{app:impl}
\begin{itemize}
  \item \textbf{Split.} The inner folds used to choose the settings are built like the outer ones, on the training
        outlets, and the settings are chosen by the Spearman correlation per outlet. In training every article counts
        once, so outlets with more articles weigh more. The language is detected with langdetect with a fixed seed.
  \item \textbf{TF-IDF.} Word uni- and bi-grams, at most 100,000, with minimum document frequency 5 and sublinear term
        frequency, fitted on the training articles; ridge penalty chosen from $\{10^{-3}, \ldots, 1\}$.
  \item \textbf{Frozen models.} The ridge regression reads the standardized vectors, with the penalty chosen from $\{10,
        \ldots, 10^6\}$; the two small frozen models often choose the lowest value, but extending the grid to 0.01
        changes their $\rho$ by less than 0.001. Qwen3-Embedding-8B and 0.6B read the article after the instruction
        \emph{``Represent this news article to judge the reliability and journalistic standards of the outlet that
        published it''}; the vector is the L2-normalized state of the last token (4,096 and 1,024 dimensions).
        Mistral-Small-24B (base) is loaded in 4-bit NF4 quantization; the vector is the mean of the hidden states of a
        layer chosen on the inner folds between layer 27 and the last one (40), and the inner folds choose layer 27 in
        every fold. The frozen ModernBERT-base is represented by the mean of its last hidden states.
  \item \textbf{Fine-tuned ModernBERT.} Squared error of $s/100$, AdamW, learning rate $3 \cdot 10^{-5}$, weight decay
        0.01, a linear schedule with 6\% warm-up, an effective batch of 32 articles, 5 epochs, mixed precision and one
        seed; the model of the last epoch is kept. Its training loss is $0.0019$ on $s/100$, an error of about 4 points,
        against 20.4 on the test articles.
  \item \textbf{Zero-shot model.} Decoding is greedy, the first integer of the answer is the prediction, and the 24
        answers without a number get the mean training score.
  \item \textbf{Length and compute.} The models read up to 4,096 tokens, more than the longest article (median about 760
        tokens). On one NVIDIA V100 (32 GB), the frozen vectors take about 2 hours for Qwen3-Embedding-8B and 5 for
        Mistral-Small-24B, fine-tuning ModernBERT takes 1.5 hours per fold, and each ridge regression a few seconds.
\end{itemize}

\section{The zero-shot prompt}\label{app:prompt}
The system message and the user message given to Mistral-Small-24B-Instruct, through its chat template:

\begin{quote}\small
You are an expert media analyst. NewsGuard rates news and information websites with a score from 0 to 100, the weighted sum of nine journalistic criteria: the site does not repeatedly publish false content (22 points), gathers and presents information responsibly (18), regularly corrects or clarifies errors (12.5), handles the difference between news and opinion responsibly (12.5), avoids deceptive headlines (10), discloses ownership and financing (7.5), clearly labels advertising (7.5), reveals who is in charge, including possible conflicts of interest (5), and provides the names of content creators with contact or biographical information (5). Scores below 60 mean that the site generally fails to meet basic standards of credibility and transparency.

\smallskip\noindent Below is one article published by a news website. The names of news outlets have been replaced by \texttt{[OUTLET]}, the names of journalists by \texttt{[AUTHOR]}, links by \texttt{[URL]} and e-mail addresses by \texttt{[EMAIL]}.

\smallskip\noindent Estimate the NewsGuard score (0-100) of the website that published this article. Answer with a single integer and nothing else.

\smallskip\noindent Article:
\emph{<the anonymized article>}
\end{quote}

\section{Robustness checks}\label{app:robust}
Table~\ref{tab:robust} reports three checks on the MAE of Table~\ref{M-tab:main}, on the same predictions:
\begin{itemize}
  \item \textbf{Recalibrated:} the outlet scores are re-mapped by a straight line fitted in every outer fold on the
        training outlets only, with predictions from three inner folds (for the zero-shot model, on its answers for the
        training outlets). This removes most of the compression and lowers every error. Mistral-Small-24B stays better
        than TF-IDF ($-1.5$ [$-2.7$, $-0.4$]), while the advantage of Qwen3-Embedding-8B is no longer significant
        ($-0.9$ [$-2.0$, 0.2]).
  \item \textbf{Re-weighted:} the MAE per NewsGuard level is weighted by the share of each level among the political
        outlets that could not be collected (43\% in the lowest level, against 15\% in the dataset) or among all the
        political outlets sent to the scraper. This estimates the error on a population of unrated outlets, about 20
        points for Mistral-Small-24B.
  \item \textbf{Without brand names:} the 16 outlets whose brand name of two or more words, as recorded by NewsGuard, is
        still in the anonymized text of at least one of their articles are left out.
\end{itemize}

\paragraph{Name only.} For each outlet, the frozen vector of its domain name alone, such as \texttt{example.com},
computed as for the articles, goes into a ridge regression over the 226 outlets with the same folds and inner choice of
the penalty. This measures how much of the score the frozen models know from the name of an outlet, without reading its
articles. Mistral-Small-24B reaches $\rho = 0.39 [0.27, 0.49]$ and an MAE of 20.4 [18.5, 22.4], and Qwen3-Embedding-8B
$\rho = 0.35 [0.23, 0.45]$ and 21.7 [19.8, 23.7]: about as much as the orientation alone ($\rho = 0.36$, MAE 18.5) and
far below what the same models reach from the articles (0.735 and 0.711).

\begin{table}[h]
\centering\small
\caption{Robustness checks. Each row is a check and each column a model; the cells give the MAE per outlet and, in
parentheses, the difference from TF-IDF in the same row (negative: better than TF-IDF).}\label{tab:robust}
\resizebox{\linewidth}{!}{%
\begin{tabular}{lccccc}
\toprule
& TF-IDF & Qwen3-Embedding-8B & Mistral-Small-24B & zero-shot & no text (median) \\
\midrule
as in Table~\ref{M-tab:main} & 17.8 & 16.3 ($-$1.5) & 15.8 ($-$2.0) & 24.4 (+6.6) & 22.1 \\
recalibrated inside the training folds & 16.2 & 15.3 ($-$0.9) & 14.7 ($-$1.5) & 21.2 (+5.0) & -- \\
re-weighted: outlets not collected & 22.9 & 20.9 ($-$2.0) & 20.2 ($-$2.7) & 22.8 ($-$0.1) & 34.9 (+12.0) \\
re-weighted: all political outlets & 22.4 & 20.4 ($-$2.0) & 19.7 ($-$2.7) & 24.0 (+1.6) & 32.6 (+10.1) \\
without the 16 outlets with their brand name & 18.0 & 16.5 ($-$1.4) & 16.2 ($-$1.8) & 24.4 (+6.5) & -- \\
\bottomrule
\end{tabular}}
\end{table}

\section{Other splits of the outlets}\label{app:splits}
The results of Table~\ref{M-tab:main} come from a single split of the outlets into five folds. We repeated the three
ridge models with four more splits, made with the same procedure and different random seeds, keeping everything else
as in Section~\ref{M-sec:setup} (Table~\ref{tab:splits}). The models are in the same order in every split and on both
metrics. Mistral-Small-24B lowers the error of TF-IDF by 2.0--2.3 points in every split (2.2 on average) and raises
$\rho$ by 0.03--0.07, and Qwen3-Embedding-8B lowers the error by 1.5--1.9 points and raises $\rho$ by 0.01--0.04. The
error of each model varies by less than half a point across splits, and the split of the paper is the one in which
Mistral-Small-24B lowers the error of TF-IDF least.

\begin{table}[h]
\centering
\caption{Other splits of the outlets. Each row is a split (the first is the split of the paper, the last the mean and
standard deviation over the five); the cells give the MAE and $\rho$ per outlet of each model, with all
articles.}\label{tab:splits}
\small\setlength{\tabcolsep}{4pt}
\begin{tabular}{lcccccc}
\toprule
& \multicolumn{2}{c}{TF-IDF} & \multicolumn{2}{c}{Qwen3-Embedding-8B} & \multicolumn{2}{c}{Mistral-Small-24B} \\
\cmidrule(lr){2-3}\cmidrule(lr){4-5}\cmidrule(lr){6-7}
split seed & MAE & $\rho$ & MAE & $\rho$ & MAE & $\rho$ \\
\midrule
1 (paper) & 17.8 & 0.674 & 16.3 & 0.711 & 15.8 & 0.735 \\
2 & 17.9 & 0.690 & 16.3 & 0.696 & 15.7 & 0.720 \\
3 & 18.2 & 0.660 & 16.5 & 0.688 & 16.0 & 0.716 \\
4 & 18.0 & 0.673 & 16.1 & 0.713 & 15.8 & 0.742 \\
5 & 18.2 & 0.673 & 16.5 & 0.692 & 16.1 & 0.722 \\
\midrule
mean (SD) & 18.0 (0.17) & 0.674 (0.010) & 16.3 (0.15) & 0.700 (0.011) & 15.8 (0.17) & 0.727 (0.011) \\
\bottomrule
\end{tabular}
\end{table}

\section{Sister outlets}\label{app:sisters}
\paragraph{Before and after deduplication.} Before deduplication (531 outlets), 298 outlets share at least one article
with another outlet. Of the 3,280 links between them, 3,212 (98\%) join outlets of the same ownership group, and so do
98\% of the 28,156 shared copies, counted per pair of outlets. The sister outlets fall from 263 to 94 at the
exact-duplicate step. After deduplication, 59\% of the 81 outlets that still have a sister (the 61 sister outlets and
the 20 outlets kept for their groups) are closer to a sister than to any outlet of another group, with a median TF-IDF
similarity of 0.73 to the nearest sister and 0.65 to the nearest outlet of another group. Ownership groups are those
recorded by NewsGuard, which may miss informal ties.

\paragraph{The sisters protocol.} The protocol of Table~\ref{M-tab:ladder} adds to the training data of every fold the
89 outlets of the same owners that deduplication or the last cleaning filter removed and that have at least 10 articles
from 2025. On the 31 outlets with a sister, the article-level split gives errors of 11.7, 12.8 and 9.6 for TF-IDF,
Qwen3-Embedding-8B and Mistral-Small-24B, about as low as with the sisters. Sisters share their ratings: 23 of the 31
outlets have the same NewsGuard score as all their sisters, outlets of the same owner have the same rating in 74\% of
the cases, and the mean absolute difference between an outlet and its sisters is 3.8 points, against 29.6 between two
outlets of these groups taken at random.

\paragraph{Similarity and gain.} For the outlets with a sister, adding the sisters raises the cosine similarity (TF-IDF)
between a test article and its most similar training article from 0.19 to 0.60 on average, against 0.16 to 0.18 for
the other outlets. Over the 31 outlets, the larger this increase, the larger the fall of the error: the Spearman
correlation between the two is 0.60 [0.30, 0.80] for TF-IDF, and 0.25 and 0.32 (not significant) for
Qwen3-Embedding-8B and Mistral-Small-24B, whose notion of similarity is probably not lexical.

\section{An earlier, rule-based anonymization}\label{app:anon}
Before the anonymization of Section~\ref{M-sec:anon}, we ran the whole pipeline, including the fine-tuned ModernBERT,
with a simpler anonymization of six rules. They masked the signature of the outlet in the titles, e-mail addresses,
links, social media accounts, the authors of the outlet taken from the author field of the page, and the name of the
outlet, that is only the publisher and its own staff. Table~\ref{tab:anon} compares the two anonymizations on the same
outlets and articles. Masking every outlet and journalist changes the language models by at most 0.010 in $\rho$ and
0.24 points of error, never significantly with all articles. Only TF-IDF loses a little ($\Delta\rho = -0.017$
[$-0.036$, $-0.001$]), as expected from a model that uses the names of outlets and people as features. With 1, 5 and
all articles per outlet and the four metrics of Table~\ref{M-tab:main}, 8 of the 12 comparisons of TF-IDF are worse
with the stricter anonymization (95\% bootstrap interval of the difference excluding zero, without correction), against
3 of 12 for Mistral-Small-24B and none for the other two models.

\begin{table}[h]
\centering
\caption{Two anonymizations. Each row is a model; the column groups give $\rho$ and MAE per outlet with the six rules of
the first run and with the anonymization of the paper (all articles).}\label{tab:anon}
\begin{tabular}{lcccc}
\toprule
& \multicolumn{2}{c}{six rules} & \multicolumn{2}{c}{rules + GLiNER} \\
\cmidrule(lr){2-3}\cmidrule(lr){4-5}
model & $\rho$ & MAE & $\rho$ & MAE \\
\midrule
TF-IDF & 0.691 & 17.7 & 0.674 & 17.8 \\
Qwen3-Embedding-8B & 0.722 & 16.1 & 0.711 & 16.3 \\
Mistral-Small-24B & 0.735 & 15.7 & 0.735 & 15.8 \\
ModernBERT & 0.605 & 17.7 & 0.610 & 17.5 \\
\bottomrule
\end{tabular}
\end{table}

\section{Boilerplate}\label{app:reading}
Many outlets open or close their articles with the same lines, such as invitations to subscribe to a newsletter,
donation appeals, disclaimers and image captions. We call boilerplate a line that appears, with the digits ignored, in
at least three articles and in at least 5\% of the articles of the same outlet. There are 1,040 such lines in 171 of the
226 outlets; they make up 3.0\% of the words and appear in 49\% of the articles.

\paragraph{TF-IDF retrained.} Without boilerplate, with the same folds and the same inner choice of the penalty, TF-IDF
has an MAE of 17.50 and $\rho = 0.677$, against 17.82 and 0.674 in Table~\ref{M-tab:main} (difference in MAE $-0.32$
[$-0.74$, 0.10]). Removing the dates as well, which can tell the model the week of a news story, gives 17.37 and 0.683
($-0.45$ [$-0.87$, $-0.03$]).

\paragraph{Saved models.} We drew 1,000 articles at random among those with boilerplate (8.7\% of their words on
average) and compared the predictions of the saved fold models, which did not see the outlet of the article, on the
original text and on the text without boilerplate (Table~\ref{tab:reading}). Without these lines the predictions fall
by 0.7--1.5 points on average and the error on these articles grows by 0.5--1.1 points, since the models were trained
with them. The two predictions remain highly correlated, most of all for Mistral-Small-24B.

\paragraph{Mistral-Small-24B token by token.} The ridge head reads the standardized mean of the $T$ token states
$\mathbf{h}_t$ of layer 27, so the prediction is $c + \frac{1}{T}\sum_{t=1}^{T}\mathbf{h}_t^\top\tilde{\mathbf{w}}$,
where $\tilde{\mathbf{w}}$ holds the ridge weights divided by the standard deviations of the features and $c$ collects
the intercept and the means. Each term of the sum is the contribution of one token. On the first 300 articles of the
sample, this decomposition matches the saved predictions within 0.04 points. The boilerplate tokens are 8.2\% of the
tokens, contribute 1.9 points to a prediction on average in absolute value (the predictions have a standard deviation of
17.5 points), and account for 12\% of the variance of the predictions (the covariance of their contribution with the
prediction divided by the variance of the prediction). They weigh slightly more than their share of the tokens, but most
of the prediction comes from the rest of the text.

\begin{table}[h]
\centering
\caption{Boilerplate. Each row is a model, on 1,000 articles that contain boilerplate; the columns give the article MAE
of the saved fold models with and without the boilerplate, the mean absolute change of the prediction, and the Pearson
correlation between the two predictions.}\label{tab:reading}
\begin{tabular}{lcccc}
\toprule
model & MAE, original & MAE, without boilerplate & mean change & $r$ \\
\midrule
TF-IDF & 19.3 & 20.4 & 3.3 & 0.954 \\
Qwen3-Embedding-8B & 18.7 & 19.5 & 4.8 & 0.936 \\
Mistral-Small-24B & 17.0 & 17.5 & 2.0 & 0.982 \\
\bottomrule
\end{tabular}
\end{table}

\section{Errors by level and orientation in detail}\label{app:errors}
\paragraph{By level.} In the 0--39 band the mean signed error is $+31$ points for Mistral-Small-24B and $+36$ for TF-IDF;
for score 100 it is $-19$ and $-25$, and in the 60--74 band $-3$ and $-2$. With Mistral-Small-24B the absolute error is
31 points in the 0--39 band, 19 for score 100 and 9 in the 60--74 band. For the trained models the signed error
decreases with the true score with a slope between $-0.63$ (Mistral-Small-24B) and $-0.81$ (Qwen3-Embedding-0.6B), so
one point of NewsGuard score moves the prediction by 0.19--0.37 points on average, and that of the zero-shot model by
0.09. The slope of the prediction on the true score is the same per article and per outlet (about 0.37 for
Mistral-Small-24B), so the aggregation keeps the compression.

\paragraph{By orientation.} Within each orientation Mistral-Small-24B reaches $\rho = 0.62$ [0.44, 0.75] among the 99
right-leaning outlets, 0.63 [0.44, 0.78] among the 59 outlets without orientation and 0.38 [0.15, 0.57] among the 68
left-leaning ones; TF-IDF and ModernBERT are lower in every group. The raw signed errors differ by orientation
(Kruskal--Wallis $p < 0.001$ for every model), mostly because the right-leaning outlets are more often in the low
levels, which are overestimated. At equal score, the gap between left- and right-leaning outlets of Mistral-Small-24B
(9.6 points [6.7, 12.3] with a linear control) becomes 8.8 with a spline of the score and 8.7 with the level as a fixed
effect. A predictor that shrinks each outlet towards the mean of its own orientation, with the slope of
Mistral-Small-24B, would have a gap of 12.6 points [9.9, 14.6], while shrinking every outlet towards the overall mean
gives no gap. At equal prediction, the true scores of left- and right-leaning outlets differ by 0.8 points [$-7.5$, 8.9]
for Mistral-Small-24B, 4.8 for Qwen3-Embedding-8B and 7.1 for TF-IDF, none significantly, and significantly by 12.8 for
the fine-tuned ModernBERT and by 24.8 for the zero-shot model.

\section{What else was tried on these outlets}\label{app:tried}
The rules of the dataset, the split and the models of the paper were fixed after earlier rounds of experiments on the
same 226 outlets, whose test results we saw. We list these experiments so that the comparisons of the paper can be read
accordingly; the ridge penalty and the Mistral layer are the only choices made inside the training folds.
\begin{itemize}
  \item We classified the five NewsGuard levels instead of regressing the score, with TF-IDF, the frozen models and
        ModernBERT.
  \item We tried other models. Gemma-27B was tested as a frozen encoder and dropped, because its half-precision outputs
        overflowed on our GPU, and layer 40 of Mistral-Small-24B is now one of the two layers chosen on the inner folds.
        DeBERTa was tested as a fine-tuned encoder in an earlier round on the exact score with three splits of the
        outlets, and the hyper-parameters of ModernBERT (5 epochs, learning rate $3 \cdot 10^{-5}$) come from an earlier
        round on these outlets.
  \item We trained fourteen other regressors on the frozen vectors of Qwen3-Embedding-8B and Mistral-Small-24B, including
        kernel ridge, a small neural network, attention over the articles of an outlet, models at the outlet level,
        weights for outlets and rare scores, and mixup. None was much better than the ridge (the best, the small neural
        network, lowered the error of Mistral-Small-24B by about one point), and the advantage over TF-IDF did not depend
        on the regressor.
  \item We re-mapped the predictions to the quantiles of the training scores, which lowers the MAE, and represented the
        outlets by groups of their articles.
  \item We used the rule-based anonymization of Section~\ref{app:anon}.
\end{itemize}
The stability of the results over five splits of the outlets (Section~\ref{app:splits}) limits how much these choices
can have adapted to the folds.

\section{Released resources and the demonstration tool}\label{app:tool}
\begin{table}[t]
\centering\small
\caption{Released resources. Each row is a released resource, with its content, the repository that holds it and its
licence. \emph{Study}: \url{https://github.com/bianchi-john/PublisherClassification}; \emph{tool}:
\url{https://github.com/bianchi-john/publisher-evaluation}\todo{publish the tool repository and check the
link}. No released file contains NewsGuard scores, NewsGuard levels or article text.}\label{tab:release}
\begin{tabular}{p{0.15\linewidth}p{0.50\linewidth}p{0.11\linewidth}p{0.12\linewidth}}
\toprule
resource & content & repository & licence \\
\midrule
pipeline code & every step of Sections~\ref{M-sec:data}--\ref{M-sec:experiments}: scraping, cleaning, anonymization,
  split, models, evaluation, tables and figures & study & Apache-2.0 \\
models & the four main models of Table~\ref{M-tab:main} (TF-IDF, Qwen3-Embedding-8B, Mistral-Small-24B,
  fine-tuned ModernBERT), each in its five fold versions (one per held-out fold of outlets):
  TF-IDF and the ridge heads as arrays, ModernBERT as ONNX & tool & Apache-2.0 \\
predictions & for each of the 25,524 articles: id, URL, outlet domain, fold and the out-of-fold score of the four
  models & tool & CC BY 4.0 \\
split & the 226 outlet domains and their fold & tool & CC BY 4.0 \\
tool & \emph{Publisher Evaluation}: scores new articles with the released models (below) & tool & Apache-2.0 \\
\bottomrule
\end{tabular}
\end{table}

The released models derive from base models under the Apache-2.0 licence: the fine-tuned encoder from ModernBERT-base
and the two frozen representations from Qwen3-Embedding-8B and Mistral-Small-24B-Base-2501, whose weights are not
redistributed but downloaded from Hugging Face by the tool. The anonymization uses the large GLiNER model of
Section~\ref{M-sec:anon} (Apache-2.0), also downloaded on first use.

\paragraph{The tool.} \emph{Publisher Evaluation} runs on the user's machine, with a web interface and a command line
(Figure~\ref{fig:tool}). Given a URL, or a pasted text with the domain of its outlet, it extracts the text as in
Section~\ref{M-sec:cleaning} and refuses what lies outside the training data instead of extrapolating: pages that are
not a single article and texts that are not in English or not between 100 and 1,500 words. The tool anonymizes the
text as in Section~\ref{M-sec:anon} and shows it with the placeholders highlighted and the list of the names replaced by
the recognizer, so the user sees exactly what the models read. Each model chosen by the user then predicts the score of
the outlet from this one article. The score of an outlet is the mean over its stored articles, shown with the standard
deviation and a bootstrap interval of this mean, which describe the spread of the article predictions and not the error
of the outlet score (about 16 points on average, Table~\ref{M-tab:main}). The released predictions of the 226 outlets of
the paper can be browsed, compared across models and exported. Every number is presented as a model prediction and
never as a rating, since one article is a weak signal (Section~\ref{M-sec:stage1}). The tool is meant for researchers
and raters, and not for publishing predictions about named outlets.

\begin{figure}[t]
\centering
\includegraphics[width=0.92\linewidth]{figures/tool.pdf}
\caption{The Evaluate page of \emph{Publisher Evaluation}, on an invented test article and with the two models that run
on a CPU. Top: the score predicted for the outlet by each model (for a new outlet, the mean of the five fold models).
Bottom: the text the models read, with the placeholders highlighted and the names replaced by the
recognizer.}\label{fig:tool}
\end{figure}

\paragraph{Model formats.} The models are shipped in formats that execute no code from the model files: TF-IDF and the
ridge heads as plain arrays rebuilt by the tool, and ModernBERT as an ONNX graph in half precision, as in our
experiments. Against the predictions of the paper, the TF-IDF scores are identical and the ridge heads differ by less
than $10^{-4}$ points on all the articles, while ModernBERT differs by 0.05 points on average (at most 1 point) on 1,500
checked articles; the anonymization of the tool reproduces that of the paper exactly on a sample of 40 articles. The two
frozen language models need a GPU with about 18 GB of memory. Without a GPU, TF-IDF and ModernBERT score new articles on
a CPU (about 1 s per article and fold for ModernBERT), and the released predictions of all four models remain
available.

\paragraph{Leakage guard.} The model of fold $k$ was trained on the outlets of the other four folds. For an article of
one of the 226 outlets, the tool therefore uses only the model of the fold that held that outlet out, so no outlet is
scored by a model that saw it in training; for a new outlet the score is the mean of the five fold models. An article of
the paper that is evaluated again keeps its identifier and is never counted twice.

\paragraph{Custom models.} Users can add their own regressors of the 0--100 score as ONNX graphs, reading either the
anonymized text or its tokens. A custom model is used only after strict checks: a valid graph with only the operators
built into ONNX and onnxruntime, exactly the declared inputs and a single score output, and a probe run on 20 synthetic
anonymized articles, in which the scores must be finite, within the scale, not all equal, and the same when scored again
and when scored in a batch. The manifest of a custom model declares which outlets it was trained on, and the same
leakage guard applies, so a model that declares nothing is never used on the 226 outlets. Its predictions are stored
with the hash of its files and never pooled with those of the models of the paper.

\end{document}
